\section{Finite geometry from the same two fixed flights}
\label{sec:two-field-finite}
The two-field inverse has a finite realization which keeps both the
flight length and the two flight directions fixed.  Its statistical
step estimates three geometric supports.  Collision locations are not
observed: every point used in a sampled hull is a commanded nominal
center at which the recorded bit was one.  Positivity permits these
same observations to serve all support directions.

Throughout this section $s=6+\beta$, the obstacles have the bounds
\eqref{eq:two-field-geometric-priors}, and the unknown footprint and density obey
\begin{equation}\label{eq:two-field-finite-priors}
 \begin{gathered}
 \diam A\le\Delta<d_0,\qquad
 \|p_A^\circ\|_{C^{6,\beta}}\le M_A,\\
 0<r_A^-\le p_A^\circ+(p_A^\circ)''\le r_A^+,
 \qquad j(z)\ge b_0\dist(z,\partial A)^\gamma
       \quad\hbox{for a.e. }z\in A,
 \qquad \gamma\ge0.
 \end{gathered}
\end{equation}
No density upper bound, continuity modulus, or density evaluations are
supplied.  Fix a dyadic $t>0$ with $2t+\Delta<d_0$, put $a=te_1$,
and use only the two ordinary forward commands $a$ and $-a$.
The unknown law is identical in every attempt.  Write
$P_C=C+(-A)$, $a_0=s(A)$, and $A_0=A-a_0$.
The width of $A$ need not be smaller than an obstacle width.

\subsection{Two-direction boundary acquisition}
The prefix inverse in Lemma~\ref{lem:two-field-prefix} already provides
a finite occupation query.  Put
\[
 K=1+\left\lceil (D_0+\Delta)/t\right\rceil.
\]
At a target $x$ its input consists of $2K$ mean values, with nominal
centers on $x+\{0,a,\ldots,Ka\}$.  If every mean has error at most
$\xi$, the computed prefix maximum has error at most $2K\xi$.
Hoeffding's inequality therefore estimates $v(x)$ to error $\eta$
with at most
\begin{equation}\label{eq:two-field-occupation-cost}
 C K^3\eta^{-2}\log\frac{CK}{\varepsilon}
\end{equation}
attempted bits and conditional failure probability at most
$\varepsilon$.  Constants and integer ceilings are understood.
This bound is uniform over the unknown law, and its aperture is fixed.
The coarse threshold uses the known boundary-mass lower bound.
Fix a small prior-dependent depth $d_*>0$ and put
$\lambda=c d_*^{\gamma+3/2}$, using the uniform cap-mass constant
in Lemma~\ref{H-lem:cap-mass}.  Estimate $v$ to error
$\lambda/4$ and retain a grid node when its estimate is at least
$\lambda/2$.  Nodes outside all $P_C$ are discarded, and every
node in $(P_C)_{-d_*}$ is retained.  This is the boundary-layer label
needed in the proof of Lemma~\ref{H-lem:coarse}; its fixed grid,
clustering, interior-ball construction and isolated radial brackets
therefore apply to $P_C$.  The cap-mass estimate is uniform under
\eqref{eq:two-field-finite-priors}, so this construction does not
require knowledge of $A$.

The following refinement uses only the two prescribed directions, even
where they are tangent to the expanded boundary.
\begin{lemma}[A two-direction rare boundary test]
\label{lem:two-field-rare}
There is a fixed collar width $d_c>0$, determined by the stated priors,
with the following property.  Fixed-accuracy occupation data supply
radial brackets for each protected $P_C$ and sufficiently accurate
normals throughout their fixed neighborhoods.  At a target $y$ in such
a neighborhood and an accuracy $0<e<d_c$, a test using only commands
$a,-a$ returns zero surely if $y\notin P_C$, and returns one with
conditional probability at least $1-\varepsilon$ if
$y\in(P_C)_{-e}$.  Its output in the intervening inner layer is
unrestricted.  The test uses at most two nominal centers and
\begin{equation}\label{eq:two-field-rare-cost}
 C e^{-(\gamma+3/2)}\log\frac C\varepsilon
\end{equation}
attempted bits.
\end{lemma}
\begin{proof}
Let $R$ be a fixed upper bound for the curvature radii of $P_C$,
increased if necessary to be positive and larger than $t$.
At a boundary point $y_0$ with outward normal $n$ one has
\begin{equation}\label{eq:two-field-outer-disk}
                 P_C\subset\overline B_R(y_0-Rn).
\end{equation}
Indeed, in support coordinates with the normal angle of $n$ as zero,
\[
 h_{P_C}(\phi)-y_0\cdot n_\phi
   =\int_0^\phi r_{P_C}(u)\sin(\phi-u)\dd u
   \le R(1-\cos\phi),\qquad 0\le\phi\le\pi.
\]
The reverse angular interval proves the same inequality for
$-\pi\le\phi\le0$.  These are exactly the support inequalities
of the disk in \eqref{eq:two-field-outer-disk}.

Choose the fixed collar so that
\[
 d_c<\min\{R/2,t^2/(32R),d_0-\Delta-t\},
\]
and below a fixed interior rolling radius of $P_C$.
Write a collar point uniquely as $y=y_0-dn$, where $|d|\le d_c$.
The coarse construction in Lemma~\ref{H-lem:rare-coarse-normals}, with
its occupation queries replaced by \eqref{eq:two-field-occupation-cost},
can give a rational $\widehat n$ satisfying
$|\widehat n-n|\le\eta_n$ uniformly on each bracket neighborhood,
where a fixed rational $\eta_n>0$ is chosen with
$6\eta_n\le t/(8R)$.  That proof uses only the uniform support,
curvature, and inball bounds of $P_C$: its finite radial interpolation
may be refined to any fixed accuracy before the fine experiment.

Set $E_1=\{e_1,-e_1\}$ and
\[
 V(\widehat n)=\{v\in E_1:\widehat n\cdot v
                  \ge\max_{w\in E_1}\widehat n\cdot w-4\eta_n\}.
\]
Every true maximizer of $n\cdot v$ is in this set.  Every candidate
satisfies $n\cdot v\ge-6\eta_n\ge-t/(8R)$.
For each candidate make ordinary forward attempts from nominal center
$y+tv$ with displacement $-tv$.  Mark the candidate if at least one
bit is one, and return the conjunction of the candidate marks.

Suppose $0<d\le d_c$.  For every candidate,
\begin{align*}
 |y+tv-(y_0-Rn)|^2-R^2
 &=t^2+2t(R-d)n\cdot v-2Rd+d^2\\
 &\ge t^2-t^2/4-t^2/16>0.
\end{align*}
Thus $y+tv$ is outside $P_C$ by
\eqref{eq:two-field-outer-disk}, so every physical start
$y+tv+z$, $z\in A$, is outside $C$.  On the event $y+Z\in C$
the attempted flight ends in $C$ and its bit is one.  Its probability
is at least $c e^{\gamma+3/2}$ if $d\ge e$, by
Lemma~\ref{H-lem:cap-mass}.

Other obstacles cannot occur in these flights.  At $y_0=c_0-z_0$,
the points $c_0\in C$ and $z_0\in A$ have normals $n$ and $-n$.
Every point of a candidate flight has distance at most
$\Delta+t+d_c<d_0$ from $c_0$.
If $d<0$, choose a true maximizing $v_*$.  Since
$n\cdot v_*\ge0$ and $n\cdot(z-z_0)\ge0$ for every $z\in A$,
\[
 n\cdot(y+tv_*+z-u tv_*)
   \ge h_C(n)-d+t(1-u)n\cdot v_*>h_C(n),\quad 0\le u\le1.
\]
This candidate always misses.  Its mark, and hence the conjunction,
is zero surely.  For inner depth at least $e$, a batch of
$C e^{-(\gamma+3/2)}\log(C/\varepsilon)$ attempts at each of the
at most two centers marks both candidates except with probability
$\varepsilon$.  Fresh attempts make the same assertion conditional
on all prior data.  No actual start or direction realization has
been observed.
\end{proof}

Apply this test in the relaxed radial bisection of
Lemma~\ref{H-lem:bisection}.  With radial angular spacing
$h\asymp\epsilon^{1/s}$ and inner layer $e\asymp\epsilon$, the
interpolation of Lemma~\ref{H-lem:radial} recovers all complete protected
expanded bodies with
\begin{equation}\label{eq:two-field-expanded-errors}
 \|\widehat h_{P_C}-h_{P_C}\|_\infty\le C\epsilon,
 \qquad
 \|\widehat h_{P_C}-h_{P_C}\|_{C^2}
                 \le C\epsilon^{(s-2)/s}.
\end{equation}
The attempted-bit cost, including coarse acquisition, is
\begin{equation}\label{eq:two-field-expanded-cost}
 C\epsilon^{-(\gamma+3/2+1/s)}
       \log\frac C\epsilon\log\frac C{\epsilon\delta}.
\end{equation}
These steps construct the expanded bodies before footprint subtraction.
Their proof and conditional confidence allocation are exactly the
radial construction just cited; its label premise is supplied by
Lemma~\ref{lem:two-field-rare}.  Rounding a target to a dyadic grid
of mesh $O(\epsilon)$ changes only the indeterminate geometric layer.
The fixed opposite vectors are exact.

\subsection{Convex hulls of positive nominal records}
Write $S_C^\pm=E_{\pm,C}$ for the closed strips in
\eqref{eq:two-field-strips}, and put
\[
                       Q_C^\pm=\operatorname{conv}(S_C^\pm)+(-A).
\]
Their support functions are those in
Lemma~\ref{lem:two-field-supports}.  The positive response itself may
have a nonconvex support.  The following sampling statement concerns
its convex hull and does not use convexity of that positive set.

\begin{lemma}[Uniform caps and shared hull acquisition]
\label{lem:two-field-hulls}
For each complete protected obstacle and each sign, the convex body
$Q_C^\pm$ can be estimated from the same two fixed commands by an
inscribed rational polygon $\widehat Q_C^\pm$ with Hausdorff error
at most $C\epsilon$, simultaneously with probability at least
$1-\delta$, using
\begin{equation}\label{eq:two-field-hull-cost}
 N\le C\epsilon^{-(\gamma+9/2)}
                         \log\frac C{\epsilon\delta}
\end{equation}
attempted bits.  A dyadic nominal grid of mesh comparable to a
sufficiently small multiple of $\epsilon^{\gamma+9/2}$ suffices.
The same positive records are used for all support directions.
\end{lemma}
\begin{proof}
First fix one strip $S=S_C^+$, the other sign being identical after
reversing $e_1$.  Its parameterization is
\[
 y=c_C(\phi)-u e_1,\qquad
          n_\phi\cdot e_1\le0,\quad 0\le u\le t.
\]
On the incoming open half-circle this parameterization is one-to-one
up to null boundary sets and has Jacobian
$r_C(\phi)|n_\phi\cdot e_1|$.  These assertions also follow by
using the transverse coordinate of each incoming boundary point.
There are uniform $c,r_0>0$ such that
\begin{equation}\label{eq:two-field-strip-ball}
             \area(S\cap B(y,r))\ge c r^3,
                     \qquad y\in S,\quad 0<r<r_0.
\end{equation}
To check the endpoint case as well, represent $y$ by $(\phi_0,u_0)$.
Choose an interval of $u$ of length $c r$ in $[0,t]$ within distance
$Ccr$ of $u_0$.  Choose an angular interval of length $c r$ within
the incoming half-circle and distance $Ccr$ of $\phi_0$, displaced
inward if $\phi_0$ is an endpoint.  It can be chosen so that
$|n_\phi\cdot e_1|\ge c' r$ on this interval.  The upper curvature
radius bounds $|c_C(\phi)-c_C(\phi_0)|$ by $Ccr$.
For a sufficiently small fixed $c$, this parameter rectangle maps
inside $B(y,r)$.  Its Jacobian, the lower curvature radius, and its
two side lengths give \eqref{eq:two-field-strip-ball}.  Take $r_0$
smaller than a fixed multiple of $t$ and of the geometric radii.

For any unit vector $u$, choose a maximizer $y_u\in S$ of
$u\cdot y$.  Equation \eqref{eq:two-field-strip-ball}, with
$r$ a small multiple of $\epsilon$, gives a subset of the strip of
area at least $c\epsilon^3$ whose support deficit in direction $u$
is at most $\epsilon/3$.  The footprint cap satisfies
\begin{equation}\label{eq:two-field-footprint-cap}
 \int_{\{z\in A: h_A(-u)+u\cdot z\le\epsilon/3\}}j(z)\dd z
                       \ge c\epsilon^{\gamma+3/2}.
\end{equation}
Indeed an interior tangent disk of fixed radius in $A$ contains a
rectangle of tangential width $c\sqrt\epsilon$ and normal depth
between $c_1\epsilon$ and $c_2\epsilon$ beneath its contact point.
Every point of a slightly smaller rectangle has distance at least
$c_3\epsilon$ from $\partial A$.  Integration of the stated lower
density bound on this rectangle proves
\eqref{eq:two-field-footprint-cap}.

Let $W$ be a fixed rational rectangle containing the strip differences
needed for all protected components and a collar.  Uniformly choosing
a continuous nominal center $X$ from $W$ and making the positive
command gives
\begin{align}
 &\Prob\{X+Z\in S_C^+,\ X\in\operatorname{cap}(Q_C^+,u,\epsilon)\}
       \nonumber\\
 &\quad=|W|^{-1}\int_Aj(z)
       \int_{S_C^+}\one_{\operatorname{cap}(Q_C^+,u,\epsilon)}(y-z)
                                         \dd y\dd z
       \ge c\epsilon^{\gamma+9/2}.
       \label{eq:two-field-positive-cap}
\end{align}
The cap consists of points $x\in Q_C^+$ satisfying
$u\cdot x\ge h_{Q_C^+}(u)-\epsilon$.  The event displayed uses
only the indicated strip, and every such event gives a positive bit
belonging to $C$, up to null boundary events.  Its obstacle label is
used only in this probability calculation.  The finite assignment
rule below recovers that label from coarse expanded-body data.
The two cap deficits add under Minkowski addition, so
the product of the preceding strip and footprint subsets proves the
inequality.

Replace $X$ by the centers of a uniform finite square grid in $W$.
Conditioning first on $Z=z$, the integrand is the indicator of a
translated strip intersected with a half-plane and $W$.  Each strip
is the difference of a convex swept body and a convex body.  The
perimeters of these bodies, and the lengths of the added straight
cuts in $W$, have a common bound.  Grid cells meeting those
boundaries have total area at most $C\ell$ for mesh $\ell$.
All other cells give the exact indicator integral.  This error bound
is uniform in $z$; integration against the probability density uses
only its total mass.  Consequently the cap probability on the finite
grid differs by at most $C\ell$ from its continuous value.
Choose a dyadic $\ell$ comparable to a sufficiently small multiple
of $\epsilon^{\gamma+9/2}$.  The lower bound in
\eqref{eq:two-field-positive-cap} then remains valid with a smaller
constant.  No regularity modulus of $j$ has been used.

The fixed coarse expanded-body hulls assign records to obstacles.
More precisely choose their Hausdorff error $\sigma_c>0$ so that
$2t+4\sigma_c<d_0-\Delta$.  A positive record of the plus or minus
command belonging to $C$ lies within distance $t+\sigma_c$ of its
coarse hull $H_C$.  Every record belonging to another obstacle has
distance at least $d_0-\Delta-t-\sigma_c$ from $H_C$.
Testing distance at the threshold $t+2\sigma_c$ therefore gives a
unique correct assignment, with a fixed reserve for finite arithmetic.
The surrounding aperture and complete-component cutoff retain every
record from each required complete component.  Exterior fragments are
not used as complete components.  Alternatively, it suffices to retain
one chosen complete component for the calibration below.

For each sign take the convex hull of the assigned positive nominal
centers.  It lies in $Q_C^\pm$ surely.  A fixed-aperture support
function is uniformly Lipschitz in direction, so a net of at most
$C\epsilon^{-1}$ unit directions with spacing $c\epsilon$ suffices
to control its supremum error by $C\epsilon$.  For each net cap the
probability of receiving no assigned positive center after $N$ attempts
is at most $\exp(-cN\epsilon^{\gamma+9/2})$.  The union bound over
the finitely many components, two signs, and this net proves
\eqref{eq:two-field-hull-cost}.  The net is only an analysis device:
the same collection of known nominal centers supplies every support
value of the resulting polygon.  All solid starts and free misses
are counted in $N$ and produce no retained point.
\end{proof}

\subsection{Stable extraction of the contact chord}
The next estimate gives the quantitative step from three support
functions to smooth geometry.  It uses the isolated negative atoms in
\eqref{eq:two-field-curvature-measure}; it does not differentiate an
empirical support function pointwise.

\begin{lemma}[A stable chord correction]
\label{lem:two-field-stable-chord}
Let $C$ satisfy the fixed smoothness and curvature priors.  Let $L$
join its two contact points with normals $e_2,-e_2$, let
$\ell$ be the length of this chord, and let $\theta$ be its normal
angle chosen with $n_\theta\cdot e_1>0$.  Put
$H=h_C-h_L$.  If a computable function $\widehat H$ has
$\|\widehat H-H\|_\infty\le\epsilon$, then finite deterministic
operations produce estimates $\widehat\theta,\widehat\ell$ and a
smooth strictly convex centered body $\widehat C^\circ$ such that
\begin{align}
 |\widehat\theta-\theta|&\le C\epsilon,
             \label{eq:two-field-chord-angle-error}\\
 |\widehat\ell-\ell|&\le C\epsilon^{(s-1)/s},
             \label{eq:two-field-chord-length-error}\\
 \|p_{\widehat C^\circ}-p_C^\circ\|_{C^2}
             &\le C\epsilon^{(s-2)/s},
             \label{eq:two-field-chord-C2-error}\\
 \|p_{\widehat C^\circ}-p_C^\circ\|_\infty
             &\le C\epsilon^{(s-1)/s}.
             \label{eq:two-field-chord-C0-error}
\end{align}
All constants are uniform in the geometric class.  Numerical
integration and representation errors can be included in these bounds.
\end{lemma}
\begin{proof}
The uniform interior radius gives a known lower chord bound
$w_*>0$, while $\ell\le D_0$ and
$n_\theta\cdot e_1\ge w_*/D_0$.  The distributional identity is
\[
 (\partial_\phi^2+1)H
   =r_C(\phi)\dd\phi-\ell(\delta_\theta+\delta_{\theta+\pi}).
\]
For $0<b<\pi/4$ set
\[
 \mathcal D_b f(\alpha)
     =f(\alpha+b)+f(\alpha-b)-2\cos b\,f(\alpha).
\]
Solving the scalar differential equation on $[\alpha-b,\alpha+b]$,
or integrating twice by parts, gives
\begin{equation}\label{eq:two-field-sine-difference}
 \mathcal D_b H(\alpha)
 =\int_{-b}^b\sin(b-|u|)r_C(\alpha+u)\dd u
    -\ell\sum_{\zeta\in\{\theta,\theta+\pi\}:|\zeta-\alpha|<b}
                                   \sin(b-|\zeta-\alpha|).
\end{equation}
If the interval contains no atom this is nonnegative.  If an atom
lies within $b/16$ of its center, it is at most
$-c\ell b+Cb^2$.  Replacing $H$ by $\widehat H$ changes the
expression by at most $4\epsilon$.

Take $b_0=L\epsilon$ with a sufficiently large fixed $L$ depending
on $w_*$.  Evaluate $\mathcal D_{b_0}\widehat H$ on an angular grid
of spacing at most $b_0/8$, with numerical error at most $\epsilon$,
and retain centers where it is smaller than $-w_*b_0/4$.
For sufficiently small $\epsilon$, there is a retained center near
each atom, and every retained center is within $b_0$ of an atom.
The criterion $n_\alpha\cdot e_1\ge w_*/(2D_0)$ selects exactly
the neighborhood of $\theta$, once $\epsilon$ is small.
Choosing one such center proves
\eqref{eq:two-field-chord-angle-error}.
The length of this computational grid has no observation cost.

Fix an even smooth function $\psi$ supported in $(-1,1)$, equal to
one on $[-1/8,1/8]$, and satisfying
$\int u^k\psi(u)\dd u=0$ for $0\le k\le4$.
Such a signed function is obtained from an even central plateau by
adding three pairs of disjoint smooth bumps outside that plateau.
Choose their positive centers distinct.  The matrix of their zeroth,
second and fourth moments is invertible when the supports are small,
by continuity from its Vandermonde limit; choose one such fixed set
and solve the three linear equations.  Odd moments vanish by symmetry.
This construction is fixed and computable.

Put $b=\epsilon^{1/s}$, up to a fixed dyadic factor, and
$\psi_b(\phi)=\psi((\phi-\widehat\theta)/b)$.
For small $\epsilon$ the true atom is in its plateau and the other
atom is outside its support.  Define
\[
 \widetilde\ell
       =-\int\widehat H(\phi)(\psi_b''(\phi)+\psi_b(\phi))\dd\phi.
\]
The true integral equals
$\ell-\int r_C\psi_b$.  Since $r_C\in C^{4,\beta}$ uniformly,
Taylor expansion about $\widehat\theta$ and the vanishing moments
bound the latter integral by $Cb^{5+\beta}=Cb^{s-1}$.
The data error is at most
$C\epsilon(b^{-1}+b)$.  Thus projection of
$\widetilde\ell$ onto $[w_*,D_0]$ proves
\eqref{eq:two-field-chord-length-error}.
An integration error of $O(\epsilon/b)$ is harmless.

Let $L^\circ$ be the centered chord, so
$h_{L^\circ}(\phi)=(\ell/2)|\sin(\phi-\theta)|$.
Let $\widehat L^\circ$ have the estimated parameters and let $T_b$
be convolution with the fixed signed approximation kernel of
Section~\ref{H-sec:stationary}, with moments through degree six
vanishing, at scale $b$.  The function
\[
              \widetilde p=T_b(\widehat H+h_{\widehat L^\circ})
\]
approximates $h_C-m\cdot n_\phi$, where $m$ is the chord midpoint.
For $f_\theta(\phi)=|\sin(\phi-\theta)|$ the curvature identity
$f_\theta''+f_\theta=2(\delta_\theta+\delta_{\theta+\pi})$
gives
\[
 \|T_b f_\theta\|_{C^2}\le Cb^{-1},\qquad
 \|\partial_\theta T_b f_\theta\|_{C^2}\le Cb^{-2}.
\]
The first derivative and supremum norms in the first bound are
bounded independently of $b$.  Consequently the error caused by
estimating the chord is at most
\[
 C\{\epsilon^{(s-1)/s}b^{-1}+\epsilon b^{-2}\}
       =C\epsilon^{(s-2)/s}
\]
in $C^2$.  Smoothing the data error costs $C\epsilon b^{-2}$,
and smoothing the true $C^{6,\beta}$ support costs $Cb^{s-2}$.
The corresponding supremum error is at most $Cb^{s-1}$.
Steiner centering, a bounded linear operation on support functions,
proves \eqref{eq:two-field-chord-C2-error} and
\eqref{eq:two-field-chord-C0-error}.

For small $\epsilon$ the uniform lower curvature radius makes this
smooth centered function an actual strictly convex support.
Its third derivatives remain uniformly bounded: the three error
terms above at derivative order three are bounded by
$C(b^{s-1}b^{-2}+\epsilon b^{-3}+b^{s-3})$.
Finite evaluation and quadratic local interpolation with the fixed
smooth partition used in the retained support construction therefore
give a finite smooth representation with any reserved $O(b^{s-2})$
$C^2$ error and $O(b^{s-1})$ supremum error.  The polygonal input
supports, their known Lipschitz bounds, and the fixed smooth kernels
permit certified finite integration at the precisions just specified.
This completes all numerical operations in the assertion.
\end{proof}

\begin{theorem}[Finite reconstruction from two fixed opposite fields]
\label{thm:two-field-finite-geometry}
Fix the priors \eqref{eq:two-field-finite-priors}, the bounded periodic
class $\mathcal B_\eta$, and one dyadic $t>0$ satisfying
$2t+\Delta<d_0$.  A finite controller using only the two forward
commands $te_1,-te_1$ at one unknown stationary launch law reconstructs
$A_0$ and a smooth strictly convex periodic table approximating
$\mathcal O-a_0$.  Its footprint support error in $C^2$ and its
geometric loss are at most $C\nu$, and its primitive discrete data
are correct, with probability at least $1-\delta$.

For sufficiently small $\nu>0$ and $0<\delta<1/4$, put
\begin{equation}\label{eq:two-field-finite-exponent}
                    Q_{\rm pair}=\frac{(\gamma+9/2)s}{s-2}.
\end{equation}
The attempted-bit and command counts satisfy
\begin{equation}\label{eq:two-field-finite-cost}
 N_\nu\le C\nu^{-Q_{\rm pair}}
              \log\frac C\nu\log\frac C{\nu\delta},
                     \qquad S_\nu\le J_\nu\le N_\nu.
\end{equation}
Every occurrence of a selected nominal center and its sign is included
in $J_\nu$, and every fresh preparation is included in $N_\nu$.
The two flight vectors and their nonzero length remain fixed as
$\nu\downarrow0$.  A nominal dyadic grid of mesh comparable to
\begin{equation}\label{eq:two-field-finite-mesh}
                         \nu^{Q_{\rm pair}}
\end{equation}
suffices.  The controller description is at most
$CJ_\nu\log(C/(\nu\delta))$ bits for fixed computable priors.
The theorem also retains finite approximations to the expanded bodies
$P_C$, their locations, and one pair of raw convex hulls
$\widehat Q_C^+,\widehat Q_C^-$.  In particular a rational polygon
approximating any selected complete $P_C$ to Hausdorff error $C\nu$
can be included in the output without additional observations.

The nominal aperture depends only on the fixed geometric priors.
Actual starts lie in its Minkowski sum with $A$, and flights in the
$t$-parallel body of that sum.  A coarse bound on $|a_0|$ gives an
absolute physical aperture; it does not supply $a_0$.  Computational
support fitting, movements, and realization of the coordinate grid
are distinct from the attempted-bit bound.  No sharpness of
$Q_{\rm pair}$ is asserted.
\end{theorem}
\begin{proof}
Choose $\epsilon$ comparable to a sufficiently small multiple of
$\nu^{s/(s-2)}$.  First obtain the fixed-accuracy expanded-body
hulls and all complete protected components.  The bounded
presentation supplies a complete obstacle in a prior-bounded nominal
region after any common translation.  Choose one such component $C_*$
for calibration, using the finite coarse hulls and the
complete-component cutoff.  Neither its label nor a free point was
an input.  Enlarge the fixed aperture to contain all required
$P_C$, both strip hulls for $C_*$, their assignment collars, and
all prefix-query sites.  The finite-cloud alternative below supplies
the corresponding complete aperture directly.

Recover the smooth expanded bodies by
\eqref{eq:two-field-expanded-errors}, and the two strip hulls of
$C_*$ by Lemma~\ref{lem:two-field-hulls}.  On their common accuracy
event form
\[
 \widehat H
    =2\widehat h_{P_{C_*}}
       -h_{\widehat Q_{C_*}^+}-h_{\widehat Q_{C_*}^-}
                                  +t|n_\phi\cdot e_1|.
\]
The support identities give
$\|\widehat H-(h_{C_*}-h_{L_*})\|_\infty\le C\epsilon$.
Lemma~\ref{lem:two-field-stable-chord} recovers the centered shape
$C_*^\circ$ with $C^2$ error $C\nu$.  Center the expanded support
and set
\[
  \widehat p_{-A_0}
        =\widehat p_{P_{C_*}}^\circ-
                                    p_{\widehat C_*^\circ}.
\]
This estimates $p_{-A_0}$ in $C^2$ error $C\nu$, since Steiner
centering commutes with Minkowski addition.  Its curvature is positive
for small $\nu$ by the footprint radius reserve.  Reflection gives
the required centered footprint.  Every original component is now
recovered by the same support subtraction
\begin{equation}\label{eq:two-field-finite-subtraction}
       \widehat h_{C-a_0}
                   =\widehat h_{P_C}-\widehat p_{-A_0}.
\end{equation}
The uniform obstacle curvature margin again gives smooth strictly
convex outputs.  Finite smooth representation, with its numerical
errors reserved, is supplied as in the last paragraph of
Lemma~\ref{lem:two-field-stable-chord}.
The observation error in the footprint and all reconstructed
components has support supremum order
$\epsilon^{(s-1)/s}$ and $C^2$ order
$\epsilon^{(s-2)/s}$.  The expanded components themselves retain the
stronger supremum order $\epsilon$.

Apply the retained period-locking and periodic assembly argument to
these original component estimates, in the frame $\mathcal O-a_0$.
The nonperiod-patch margin is invariant under a common translation,
as in Lemma~\ref{H-lem:registration-period-margin}.  The bounds just
obtained tend to zero uniformly, so the fixed locking thresholds are
met.  Primitive relations and orbit counts are exact on the common
event; a matched period basis and free area have error
$O(\epsilon^{(s-1)/s})$, which is $O(\nu)$.

The hull cost \eqref{eq:two-field-hull-cost} dominates the expanded
boundary cost \eqref{eq:two-field-expanded-cost}, since
$\gamma+9/2>\gamma+3/2+1/s$.  Splitting the confidence budget among
coarse acquisition, adaptive boundary queries, and the two hull
experiments proves \eqref{eq:two-field-finite-cost} after substituting
$\epsilon\asymp\nu^{s/(s-2)}$.
The nominal grid in Lemma~\ref{lem:two-field-hulls} is finer than the
$O(\epsilon)$ geometric rounding required in the boundary queries.
Choose it dyadically so that $t$ is an integer multiple of its mesh.
All shifts then preserve the grid, and
\eqref{eq:two-field-finite-mesh} follows.  Only two displacements
occur.  Every sampled grid point has a finite index of
$O(\log(C/\nu))$ bits in the fixed aperture.  Recording these indices,
signs, and batch counts gives the stated description bound.

Finally, the expanded-body estimates have uniformly bounded curvature
radii.  Sampling their smooth boundaries at sufficiently fine
numerical angular spacing and rounding those values gives inscribed
or nearby rational polygons of any reserved $O(\nu)$ Hausdorff error.
This is a finite computation on the already acquired functions and
uses no new collision records.  Their locations are those of the
measured nominal-coordinate bodies $P_C$.
\end{proof}

\begin{corollary}[A complete finite configuration]
\label{cor:two-field-finite-cloud}
Suppose instead there are between one and $n_0$ obstacles, with the
same uniform priors, and a known protected aperture contains every
expanded body $C-A$ and the fixed collars above.  The same controller
recovers $A_0$ and all of $\mathcal O-a_0$ in $C^2$ error $C\nu$
with probability at least $1-\delta$ and the bounds
\eqref{eq:two-field-finite-cost}--\eqref{eq:two-field-finite-mesh}.
No periodicity or positive nonperiod-patch margin is used.
\end{corollary}
\begin{proof}
The complete aperture supplies all coarse expanded components and a
calibration component.  The hull acquisition, chord correction, and
common support subtraction in the theorem recover every body.
Omit the period test and periodic assembly; all preceding estimates
and observation counts remain unchanged.
\end{proof}
